\documentclass[runningheads]{llncs}

\usepackage{graphicx}
\usepackage{amsmath}
\usepackage{amssymb}
\usepackage{booktabs}
\usepackage{microtype}
\usepackage{hyperref}
\usepackage{url}

\begin{document}

\title{The Zeitgeist Hypothesis: Subcommunity Power Laws and\\
Temporal Localization in a 130-Year Citation Corpus}

\author{David Goh\inst{1} \and Xiaobai Sun\inst{1}}

\authorrunning{D. Goh and X. Sun}

\institute{Department of Computer Science, University of Toronto,
Toronto, ON M5S 2E4, Canada\\
\email{daveed@cs.toronto.edu}}

\maketitle

\begin{abstract}
The global in-degree distribution of citation networks is commonly modeled as a
single power law, yet this aggregate view obscures the heterogeneous citation
dynamics of distinct research communities.  We articulate and empirically test
the \emph{Zeitgeist hypothesis}: the global citation distribution of a large
scientific corpus is a finite mixture of subcommunity distributions, each
individually scale-free and each corresponding to a temporally localized
research generation.  Using the American Physical Society (APS) 2022 corpus
(709,803 papers; 9,833,191 citations; 1893--2022), we decompose the corpus into
communities via the Leiden algorithm (446 communities, modularity $Q=0.7883$)
and apply discrete maximum-likelihood estimation with $K_{\min}$ scanning and
Kolmogorov--Smirnov testing to each large community ($n \geq 30$).  All 25
large communities pass the KS test ($p \geq 0.05$), with fitted exponents
$\gamma_c \in [2.099,\, 3.268]$ (mean $2.50 \pm 0.25$).  Furthermore, 68\% of
communities have a publication-year interquartile range under 20 years, and
community year-medians span 1950 to 2017 --- confirming temporal localization.
Our results provide the first empirical validation of the mixture structure of
citation dynamics on a 130-year longitudinal corpus.
\end{abstract}

\keywords{citation networks \and power laws \and community detection \and
temporal localization \and scientometrics}

%-----------------------------------------------------------------------
\section{Introduction}
\label{sec:intro}
%-----------------------------------------------------------------------

Academic citation networks encode the causal structure of knowledge propagation.
Their edges are directed, reflecting the irreversibility of intellectual
influence; their nodes are timestamped; and their bibliographies are
accountable, making them among the most reliable social networks available for
empirical study.  These properties have sustained decades of research into
the statistical regularities of scientific output
\cite{Price1965,Price1976,Redner1998,Newman2003}.

A central empirical finding is that citation in-degree distributions follow an
approximate power law: $P(k) \sim k^{-\gamma}$.  Preferential attachment
\cite{Price1976} provides the canonical generative explanation --- papers with
more citations attract more citations, producing a ``rich-get-richer'' dynamic.
Barab{\'a}si \cite{Barabasi2016} fit this model to the APS corpus and reported
$\gamma = 2.79$ with $K_{\min} = 49$, or $\gamma = 3.03$ with an exponential
saturation cutoff.  He also noted, however, that the pure power law fails a
formal goodness-of-fit test at the corpus level ($p < 10^{-4}$).

We argue that this global-model failure is not a statistical artifact but
a structural signal.  Consider the epistemic implausibility of a single
preferential attachment process governing a 1935 nuclear structure paper and a
2010 topological insulator paper simultaneously.  These papers inhabited
different research communities, drew on different citation pools, and attracted
citations over different decades.  The global distribution lumps them together;
the per-community distribution respects their separateness.

We therefore propose the \textbf{Zeitgeist hypothesis}: the global citation
distribution of a scientific corpus is a finite mixture of subcommunity
distributions,
\begin{equation}
  P(k) = \sum_{c} \alpha_c\, P_c(k), \qquad P_c(k) \sim k^{-\gamma_c},
  \label{eq:mixture}
\end{equation}
where each community $c$ has weight $\alpha_c > 0$ and an individually
scale-free distribution with exponent $\gamma_c$.  Two empirical predictions
follow from Eq.~\eqref{eq:mixture}: \emph{(i)}~each subcommunity should pass a
power-law goodness-of-fit test applied in isolation, and \emph{(ii)}~the
subcommunities should be temporally localized --- each corresponding to a
research generation concentrated in a recognizable historical window.

We test both predictions on the APS 2022 corpus using a three-stage pipeline:
(1) community detection via the Leiden algorithm \cite{Traag2019}; (2)
per-community power-law fitting via discrete MLE with $K_{\min}$ scanning
\cite{Clauset2009}; and (3) temporal analysis of publication-year distributions
within each community.  The results confirm the Zeitgeist hypothesis: all 25
large communities individually exhibit power-law citation behavior, and their
temporal footprints are narrow and well-separated across 130 years.

To our knowledge, this is the first empirical study to jointly test the
mixture-decomposition claim and the temporal-localization claim on a
longitudinal corpus of this scale.


%-----------------------------------------------------------------------
\section{Related Work}
\label{sec:related}
%-----------------------------------------------------------------------

\paragraph{Scale-free citation networks.}
Power-law behavior in citation distributions was documented as early as
Price \cite{Price1965,Price1976}, who also introduced the preferential
attachment mechanism.  Redner \cite{Redner1998} confirmed approximate power-law
tails in Physical Review data.  Barab{\'a}si \cite{Barabasi2016} fit the full
APS corpus and provided the benchmark $\gamma = 2.79$ against which we
calibrate our global-fit result.  Ke et al.\ \cite{Ke2023} studied the temporal
pattern of citation accumulation, identifying a ``golden age'' window after
which paper impact decays; their temporal framing is complementary to ours but
operates at the individual paper level rather than the community level.

\paragraph{Community structure in citation networks.}
Waltman and van Eck \cite{Waltman2012} developed community-based classification
systems for science from citation data, establishing that citation-derived
communities align with disciplinary boundaries.  Blondel et al.'s Louvain
method \cite{BlondelLouvain2008} and the improved Leiden algorithm
\cite{Traag2019} make large-scale community detection tractable; we use Leiden
for its guarantee of well-connected communities.  Apar{\'i}cio et al.\
\cite{Aparicio2024} studied temporal community trajectories in the KDD citation
network, the work closest in spirit to our temporal localization analysis, but
they did not test power-law behavior within communities.

\paragraph{Mixture models and phase structure.}
Castillo-Castillo et al.\ \cite{CastilloCastillo2025} recently proposed a
generative mixture model for citation dynamics, providing a formal probabilistic
framework for Eq.~\eqref{eq:mixture}.  Our contribution is empirical: we
validate the mixture claim directly on the APS corpus using non-parametric
goodness-of-fit testing rather than fitting a generative model.  Costa and
Frigori \cite{CostaFrigori2024} characterized phase transitions in citation
networks of AI literature using entropy and fractal dimension, offering a
different formalism for temporal change; their scope is narrower (one research
field) and they do not decompose into subcommunities.

\paragraph{Gap.}
No prior work jointly tests the mixture-decomposition claim
(each subcommunity is individually power-law) and the temporal-localization
claim on a 130-year longitudinal corpus.


%-----------------------------------------------------------------------
\section{Dataset}
\label{sec:data}
%-----------------------------------------------------------------------

We use the APS 2022 citation dataset, which covers all articles published in
20 American Physical Society journals from 01 July 1893 to 30 December 2022.
The dataset associates each article with a unique DOI and metadata including
journal, volume, issue, publication date, and author affiliations.

\paragraph{Graph statistics.}
We construct a directed citation graph $G = (V, E)$ where each node
$v \in V$ is a paper and each edge $(u, v) \in E$ indicates that $u$ cites $v$.
The resulting graph has:
\[
  |V| = 709{,}803 \text{ papers}, \quad
  |E| = 9{,}833{,}191 \text{ citations}.
\]
The graph has mean in-degree $\bar{k} = 13.84$, Gini coefficient $G = 0.692$,
and contains one giant weakly connected component alongside 447 small
components.  Year metadata is present for $99.3\%$ of nodes
(1893--2022, spanning 130 years).

\paragraph{Concentration.}
Citation impact is highly concentrated: the top $1\%$ of papers account for
$20.5\%$ of all citations, and the top $20\%$ account for $71.5\%$.  This Pareto
structure is consistent with preferential attachment dynamics.

\paragraph{Global degree distribution.}
Figure~\ref{fig:ccdf} shows the in-degree complementary cumulative distribution
function (CCDF) on a log-log scale.  Applying discrete MLE with $K_{\min}$
scanning over $[1, 100]$ and 500 KS bootstrap samples \cite{Clauset2009}
yields $\hat{\gamma}_{\text{global}} = 2.738$ with $K_{\min} = 96$, consistent
with Barab{\'a}si's benchmark of $\gamma = 2.79$ \cite{Barabasi2016}.  As
previously noted \cite{Barabasi2016}, the pure power law fails a formal
goodness-of-fit test at the corpus level, motivating the mixture decomposition
we develop in Section~\ref{sec:zeitgeist}.

\begin{figure}[t]
  \centering
  \includegraphics[width=0.75\linewidth]{../data/figures/fig1_indegree_ccdf.pdf}
  \caption{In-degree CCDF of the APS corpus on a log-log scale.  The fitted
    power-law tail ($\hat{\gamma} = 2.74$, $K_{\min} = 96$) is shown as a
    dashed line.  The global distribution is approximately scale-free but
    formally rejects the pure power-law model, motivating the community-level
    mixture decomposition.}
  \label{fig:ccdf}
\end{figure}


%-----------------------------------------------------------------------
\section{The Zeitgeist Hypothesis}
\label{sec:zeitgeist}
%-----------------------------------------------------------------------

\subsection{Formal Statement}

We formalize the Zeitgeist hypothesis as follows.  Let $G = (V, E)$ be a
directed citation graph with a community partition $\mathcal{C} =
\{C_1, \ldots, C_K\}$ of $V$.  The Zeitgeist hypothesis asserts:

\begin{enumerate}
  \item \textbf{Local scale-free property:} For each community $C_c$ of
        sufficient size, the in-degree distribution of the induced subgraph
        $G[C_c]$ follows a power law $P_c(k) \sim k^{-\gamma_c}$ for $k \geq
        K_{\min,c}$.
  \item \textbf{Temporal localization:} For each community $C_c$, the
        publication years of its member papers are concentrated in a bounded
        historical window; specifically, the interquartile range (IQR) of
        publication years is small relative to the corpus span.
\end{enumerate}

Both properties follow naturally if each community corresponds to a research
generation in which papers compete primarily for citations within a shared pool.
Global power-law failure then reflects the heterogeneity across generations, not
the absence of scale-free dynamics within them.

\subsection{Community Detection}

We partition the APS citation graph using the Leiden algorithm \cite{Traag2019}
at resolution $\lambda = 1.0$, applied to the full $|V| = 709{,}803$ node graph.
Leiden partitions the graph into $K = 446$ communities with modularity
$Q = 0.7883$.  Figure~\ref{fig:sizes} shows the community size distribution.

\begin{figure}[t]
  \centering
  \includegraphics[width=0.75\linewidth]{../data/figures/fig2_community_sizes.pdf}
  \caption{Community size distribution from Leiden partitioning ($K = 446$
    communities, $Q = 0.7883$).  Twenty-five communities have $n \geq 30$
    nodes, collectively covering 99.8\% of all papers (708,622 of 709,803).
    The remaining 421 communities are small fragments ($n < 30$), typical
    at this resolution.}
  \label{fig:sizes}
\end{figure}

The size distribution is extremely heavy-tailed: 25 communities have
$n \geq 30$ nodes and together cover $708{,}622 / 709{,}803 = 99.8\%$ of all
papers.  The remaining 421 communities are fragments of fewer than 30 nodes
--- expected behavior of Leiden at resolution $\lambda = 1.0$ on a graph
with a dominant giant component.  All subsequent analyses focus on the 25
large communities.

\subsection{Per-Community Power-Law Fitting}

For each large community $C_c$, we extract the in-degree sequence of the
induced subgraph $G[C_c]$ and apply the discrete MLE procedure of Clauset et
al.\ \cite{Clauset2009} with $K_{\min}$ scanning.  Specifically, we search
$K_{\min} \in [1, 100]$, select the value minimizing the KS statistic between
the empirical CCDF and the fitted power law, and report the fitted exponent
$\hat{\gamma}_c$.  We assess goodness of fit via 500 KS bootstrap samples;
a community passes if $p \geq 0.05$.

\paragraph{Results.}
All 25 large communities pass the KS test ($p \geq 0.05$), providing strong
support for the local scale-free property.  The fitted exponents span
$\hat{\gamma}_c \in [2.099,\, 3.268]$, with mean $2.500$ and standard deviation
$0.246$.  Figure~\ref{fig:gamma} shows the distribution of $\hat{\gamma}_c$
across communities.

\begin{figure}[t]
  \centering
  \includegraphics[width=0.75\linewidth]{../data/figures/fig3_gamma_histogram.pdf}
  \caption{Distribution of fitted power-law exponents $\hat{\gamma}_c$ across
    the 25 large APS communities.  All communities pass the KS goodness-of-fit
    test ($p \geq 0.05$).  The mean exponent ($\mu = 2.50$, shown as a dashed
    line) is consistent with the global fit ($\hat{\gamma} = 2.74$), but the
    range $[2.099, 3.268]$ reveals genuine inter-community variation masked by
    the aggregate.}
  \label{fig:gamma}
\end{figure}

The global exponent ($\hat{\gamma}_{\text{global}} = 2.74$) lies within the
range of per-community exponents, consistent with the mixture interpretation of
Eq.~\eqref{eq:mixture}: the global distribution is a weighted combination of
individually valid power laws, and the aggregate exponent reflects the
community-size-weighted mean.

\subsection{Temporal Localization}

We compute, for each large community $C_c$, the median publication year
$\tilde{y}_c$ and the year interquartile range $\text{IQR}_c$.  Figure~\ref{fig:timeline}
visualizes these temporal footprints as horizontal bars sorted by $\tilde{y}_c$.

\begin{figure}[t]
  \centering
  \includegraphics[width=\linewidth]{../data/figures/fig4_timeline.pdf}
  \caption{Publication-year footprint of the 25 large APS communities, sorted
    by median year.  Each bar shows the IQR; the dot marks the median.
    Community medians span from 1950 (Early Nuclear/Atomic Physics) to 2017
    (Topological Matter/Graphene and Physics Education Research).  Seventeen of
    25 communities have IQR $< 20$ years.}
  \label{fig:timeline}
\end{figure}

\paragraph{Results.}
The mean year IQR across the 25 communities is 18.4 years (median 17.0 years).
Seventeen of 25 communities ($68\%$) have IQR $< 20$ years, confirming that the
majority of research communities are temporally concentrated into windows well
under one research generation.  Community year-medians span from 1950 (Community
20, Early Nuclear/Atomic Physics) to 2017 (Community 12, Topological
Matter/Graphene; Community 21, Physics Education Research), covering the full
historical arc of modern physics.

\paragraph{Top communities.}
Table~\ref{tab:top10} reports the 10 largest communities with their fitted
parameters and physics labels.  The labels were assigned by inspecting the top
five most-cited DOIs per community and cross-referencing with APS journal
context.  The communities recapitulate the major subfields of APS physics across
the 20th and 21st centuries.

\begin{table}[t]
\caption{Top-10 communities by size.  KS: all 25 large communities pass the
  Kolmogorov--Smirnov power-law test ($p \geq 0.05$).
  $\hat{\gamma}_c$ fitted via discrete MLE with $K_{\min}$ scanning
  \cite{Clauset2009}.  yr = publication year.}
\label{tab:top10}
\centering
\small
\begin{tabular}{rrrccrl}
\toprule
cid & $n$ & $\hat{\gamma}_c$ & KS & yr med & yr IQR & Label \\
\midrule
0 & 93{,}221 & 2.48 & \checkmark & 1999 & 23\,y & Cond.\ Matter --- DFT / Electronic Structure \\
1 & 62{,}216 & 2.53 & \checkmark & 2003 & 20\,y & Cond.\ Matter --- Magnetism / Disordered Systems \\
2 & 55{,}859 & 2.63 & \checkmark & 1987 & 37\,y & Nuclear Physics \\
3 & 54{,}598 & 2.37 & \checkmark & 1997 & 35\,y & Particle Physics --- Field Theory / QCD \\
4 & 52{,}502 & 2.29 & \checkmark & 2006 & 17\,y & Mesoscopic Physics / Quantum Chaos \\
5 & 49{,}428 & 2.24 & \checkmark & 2011 & 15\,y & Quantum Information / Computing \\
6 & 47{,}156 & 2.68 & \checkmark & 1992 & 32\,y & AMO Physics / Quantum Optics \\
7 & 45{,}155 & 2.33 & \checkmark & 2010 & 17\,y & Astrophysics / Gravitational Waves \\
8 & 37{,}730 & 2.54 & \checkmark & 2007 & 16\,y & High-Temp.\ Superconductivity \\
9 & 36{,}510 & 2.32 & \checkmark & 2002 & 16\,y & Cold Atoms / BEC / Laser Cooling \\
\bottomrule
\end{tabular}
\end{table}

Together, these results validate both predictions of the Zeitgeist hypothesis:
per-community power-law behavior is universal across all 25 large communities,
and the communities are temporally localized into recognizable research
generations.


%-----------------------------------------------------------------------
\section{Discussion}
\label{sec:discussion}
%-----------------------------------------------------------------------

\paragraph{Summary.}
We have shown that the APS citation corpus, while approximately power-law in
aggregate, decomposes into 25 large communities each of which individually
passes a stringent KS goodness-of-fit test for power-law behavior.  The
fitted exponents $\hat{\gamma}_c \in [2.099, 3.268]$ span a range invisible in
the global aggregate, and 68\% of communities are temporally concentrated into
windows narrower than 20 years.  The community year-medians form a continuous
sequence from 1950 to 2017, tracing the historical succession of APS research
generations from early nuclear physics through quantum information and
topological matter.

\paragraph{What the exponent variation means.}
The range $\hat{\gamma}_c \in [2.099, 3.268]$ with $\sigma = 0.246$ is
moderate but meaningful.  A lower $\gamma_c$ corresponds to a heavier tail ---
a community in which a few papers accumulate citations far more
disproportionately than others.  Quantum Information / Computing ($\gamma_c =
2.24$, yr median 2011) and Mesoscopic Physics / Quantum Chaos ($\gamma_c = 2.29$,
yr median 2006) have the heaviest tails, suggesting winner-takes-all citation
dynamics in newer, rapidly forming communities.  By contrast, AMO Physics
($\gamma_c = 2.68$) and Nuclear Physics ($\gamma_c = 2.63$) --- both older,
more diffuse communities with broad temporal IQRs --- have lighter tails.  The
Zeitgeist is thus not only about \emph{when} a research generation dominated,
but also about \emph{how concentrated} its citation flow became.

\paragraph{Temporal localization as a research generation signature.}
The narrow year IQRs of communities like Quantum Information (IQR 15\,y) and
Topological Matter/Graphene (IQR 7\,y) reflect the rapid rise and consolidation
of new research programs: a body of foundational papers appears within a short
window, and the community then cites that shared base intensively.  The
wider IQRs of Nuclear Physics (37\,y) and Particle Physics (35\,y) reflect their
longer historical roots --- these communities span multiple decades of APS
publication and have accumulated citations across multiple research generations
within the same field.

\paragraph{Relationship to the global fit.}
The global exponent $\hat{\gamma}_{\text{global}} = 2.74$ is consistent with
Barab{\'a}si's $\gamma = 2.79$ \cite{Barabasi2016} and lies within the range of
per-community exponents.  The global model is not wrong as a descriptive
summary, but it is not a sufficient characterization: it masks $0.246$ standard
deviations of genuine inter-community variation and provides no information
about the temporal structure of citation dynamics.

\paragraph{Limitations.}
Our analysis is restricted to the APS corpus.  It is an open question whether
the Zeitgeist hypothesis holds in biomedical, social-science, or computer-science
corpora, where citation norms and community boundaries differ.  Community
detection is resolution-dependent: our choice of $\lambda = 1.0$ produces 446
communities; a different resolution would merge or split communities and could
affect the per-community KS results.  Finally, the community labels were
assigned manually from the top-cited DOIs, not from semantic text features;
communities spanning multiple subfields may be incompletely characterized.

\paragraph{Future directions.}
A natural extension is to fit an aging model $\pi_c(t)$ per community,
testing whether the temporal IQR predicts $\hat{\gamma}_c$ --- i.e., whether
younger, more temporally concentrated communities have heavier citation tails.
Cross-corpus validation on open datasets (Semantic Scholar Open Research Corpus,
bioRxiv) would test generalizability.  Semantic community labeling via LLM
summarization of abstracts would remove the manual bottleneck and enable
scalability to millions of papers.


%-----------------------------------------------------------------------
\section*{Acknowledgments}
%-----------------------------------------------------------------------

This work uses the APS 2022 citation dataset, made available by the American
Physical Society for non-commercial research purposes.

\bibliographystyle{splncs04}
\bibliography{zeitgeist_refs}

\end{document}
